\documentclass[10pt]{article}

\usepackage[a4paper,margin=1in]{geometry}
\usepackage[T1]{fontenc}
\usepackage[utf8]{inputenc}
\usepackage{lmodern}
\usepackage{microtype}
\usepackage{booktabs}
\usepackage{amsmath}
\usepackage{amssymb}
\usepackage{array}
\usepackage{fancyvrb}
\usepackage[numbers,sort&compress]{natbib}
\usepackage[hidelinks]{hyperref}

\newcommand{\pcv}{\textsf{purgedcv}}
\newcommand{\pkg}[1]{\textsf{#1}}
\DeclareRobustCommand{\code}[1]{\texttt{\detokenize{#1}}}
\newcommand{\doi}[1]{\href{https://doi.org/#1}{\nolinkurl{#1}}}

\hypersetup{
  pdftitle={Label-aware cross-validation for overlapping-horizon prediction},
  pdfauthor={Evgenii Lazarev},
  pdfsubject={Leakage-aware cross-validation for time-series and panel prediction},
  pdfkeywords={cross-validation, data leakage, time series, panel data, model validation, reproducible software}
}

\DefineVerbatimEnvironment{Code}{Verbatim}{fontsize=\small,frame=single}

\title{Label-aware cross-validation for overlapping-horizon prediction}

\author{
Evgenii Lazarev\textsuperscript{1,*}\\
\textsuperscript{1}Independent Researcher\\
\textsuperscript{*}\href{mailto:elazarev@gmail.com}{elazarev@gmail.com}\\
\href{https://orcid.org/0009-0000-1398-7842}{ORCID: 0009-0000-1398-7842}
}

\date{}

\begin{document}
\maketitle

\begin{abstract}
Cross-validation is routinely used to estimate out-of-sample performance in
statistical learning, but standard shuffled or blocked folds can be invalid when
responses are measured over future intervals. A label such as the mean demand
over the next twelve half-hours, the next-day rainfall amount, or the return
over the next twenty bars overlaps the labels of nearby rows. If overlapping
label intervals are split between training and test sets, the validation score
partly measures information reuse rather than generalization. This study
formalizes split-level conditions for leakage-aware validation in
overlapping-label time-series and panel data, and evaluates them through
\pcv, an open-source Python implementation of purging, embargoing,
walk-forward validation, group-purged folds, and combinatorial purged
cross-validation. A controlled experiment with an unpredictable target shows
that shuffled k-fold can report a mean out-of-sample $R^2$ of 0.918 while
admitting complete train/test label overlap. A full-population benchmark on Low
Carbon London smart-meter data shows a more nuanced case: the temporal leakage
gap is small but measurable, whereas the larger issue is household-level
generalization. The results show that validation choice is a scientific
measurement decision, not only a software parameter.
\end{abstract}

\noindent\textbf{Keywords:} cross-validation; data leakage; time series; panel
data; model validation; reproducible software

\section{Introduction}

Cross-validation is often treated as a neutral measurement device: choose a
splitter, fit the same estimator on each training fold, and average the test
scores. That view depends on an independence assumption that is not satisfied by
many time-indexed prediction tasks. In a forecasting or backtesting problem, row
$i$ is usually not just an instantaneous observation. Its response may be
defined by an evaluation interval that starts at the prediction time and ends
after a future horizon. Nearby rows can therefore share part of the same future
outcome window. Standard shuffled k-fold cross-validation can place one row in
the test fold and another row whose label interval overlaps it in the training
fold. The resulting score is then contaminated by information that would not be
available when the model is used prospectively.

Data leakage is a well-known source of inflated performance estimates in
machine learning \citep{kaufman2012}. It is especially damaging in scientific
applications where a model is selected or reported after many iterations, since
the optimistic score becomes part of the published evidence rather than merely a
development mistake \citep{mcdermott2021}. In financial machine learning,
\citet{lopezdeprado2018afml} proposed purging and embargoing as practical
guards against leakage from overlapping labels and serial dependence, and
Combinatorial Purged Cross-Validation (CPCV) as a way to obtain multiple
out-of-sample backtest paths. Bailey and Lopez de Prado's Probabilistic Sharpe
Ratio and Deflated Sharpe Ratio address the related problem of selection bias
after trying multiple strategies \citep{bailey2012psr,bailey2014dsr}.

The same validation problem appears outside finance. Household electricity
demand, equipment degradation, rainfall, clinical monitoring, and air-quality
forecasting all contain future-horizon labels or repeated entities. What matters
is whether any training-label window overlaps a test-label window, whether a
post-test serial-dependence buffer has been respected, and whether the
deployment target involves unseen entities rather than future observations from
already-seen entities.

This study makes three contributions. First, it states a directly checkable
interval condition for overlapping-label validation. Second, it presents \pcv,
an open Python implementation of purging, embargoing, walk-forward validation,
group-purged k-fold, and CPCV path reconstruction through the
\pkg{scikit-learn} cross-validation interface, together with leakage
diagnostics for auditing splits \citep{pedregosa2011sklearn,purgedcv2026}.
Third, it reports reproducible experiments that show both dramatic and
undramatic outcomes: a synthetic task where leakage fabricates strong skill, and
a real smart-meter benchmark where the larger issue is not temporal leakage but
household-level generalization.

\section{Results}

\subsection{Overlapping labels define a split-level invariant}

Let a supervised learning data set contain observations
\[
  z_i = (x_i, y_i, p_i, e_i, g_i), \quad i = 1,\ldots,n,
\]
where $x_i$ is the feature vector, $y_i$ is the response, $p_i$ is the
prediction time, $e_i$ is the evaluation time at which the response is fully
known, and $g_i$ is an optional group identifier such as a household, patient,
engine, or season. The label interval for row $i$ is
\[
  I_i = [p_i, e_i).
\]
The interval is half-open: a label window ending exactly where another begins
is not counted as overlapping. For a train/test split $(A, B)$,
label-overlap leakage is present when
\[
  \exists i \in A, \exists j \in B
  \quad \text{such that} \quad I_i \cap I_j \ne \emptyset .
\]
A leakage-aware split must remove such training rows before fitting the
estimator. In practice, a second guard is often needed after a test block. If the
process remains serially correlated after the test interval, training
immediately after the test block can still reuse information tied to that test
period. An embargo removes training rows whose prediction time lies inside a
post-test buffer of fixed duration or fixed fraction of the sample.

For panel data there is a separate deployment question. If the intended use is
prediction on new entities, a chronological split that mixes the same entity
across training and test sets may answer the wrong question even when no label
intervals overlap. In that case the split must also satisfy
\[
  \{g_i: i \in A\} \cap \{g_j: j \in B\} = \emptyset .
\]
Thus validation has at least three distinct requirements: interval disjointness,
post-test embargo, and group disjointness. They are not interchangeable. A fixed
integer gap may remove leakage for a single constant horizon, but it does not
express variable horizons, time-duration embargoes, CPCV test blocks, or
entity-level generalization.

\subsection{The software makes validation contracts executable}

\pcv{} implements the interval operations and splitters needed to make these
requirements executable. The package is written in Python, is MIT licensed, and
follows the \pkg{scikit-learn} splitter protocol, so the same objects can be
passed to \code{cross_val_score}, \code{GridSearchCV}, and \code{Pipeline}.
Runtime dependencies are intentionally small: \pkg{numpy}, \pkg{pandas},
\pkg{scikit-learn}, and \pkg{scipy}
\citep{harris2020numpy,mckinney2010pandas,pedregosa2011sklearn,virtanen2020scipy}.
Table~\ref{tab:api} summarizes the public components exposed by the package.

\begin{table}[htbp]
\centering
\caption{Main public API in \pcv.}
\label{tab:api}
\footnotesize
\begin{tabular}{p{0.16\textwidth}p{0.34\textwidth}p{0.42\textwidth}}
\toprule
Component & Function or class & Purpose \\
\midrule
Primitive & \code{purge} & Drop training rows whose label intervals overlap test labels \\
Primitive & \code{apply_embargo} & Drop post-test training rows inside an embargo buffer \\
Splitter & \code{WalkForwardSplit} & Expanding or rolling chronological validation \\
Splitter & \code{PurgedKFold} & Contiguous folds with label-aware purging and embargo \\
Splitter & \code{PurgedGroupKFold} & Purged folds with disjoint held-out groups \\
Splitter & \code{CombinatorialPurgedCV} & CPCV folds with multiple test-block combinations \\
Paths & \code{reconstruct_paths} & Assemble CPCV folds into backtest paths \\
Metrics & \code{probabilistic_sharpe_ratio} & Probability that skill exceeds a benchmark \\
Metrics & \code{deflated_sharpe_ratio} & Sharpe-ratio inference corrected for selection bias \\
Metrics & \code{min_track_record_length} & Minimum observations needed to establish a Sharpe ratio \\
Diagnostics & \code{assert_*} functions & Check temporal, embargo, and group-leakage invariants \\
\bottomrule
\end{tabular}
\end{table}

The implementation separates interval arithmetic from splitter orchestration.
For each fold, test label intervals are sorted and merged once, and candidate
training intervals are tested against the merged set. This avoids duplicating
boundary logic across splitters and is particularly important for CPCV, where
the test set may contain several non-adjacent blocks. In such a fold, purging
must apply to the union of test label intervals, not to the convex hull between
the first and last test block.

The diagnostic functions are deliberately independent of the package's own
splitters. They accept training indices, test indices, prediction times, and
evaluation times, and can audit a split produced by any library or by hand. This
turns the validation contract into an assertion that can be placed in tests.

\subsection{A controlled task exposes complete validation failure}

The controlled task is designed so that no feature has genuine predictive
content. Let $\epsilon_t$ be independent noise and define the response at row
$t$ as the mean of the next $H$ future noise values. The only feature is a
monotone clock. A model cannot forecast the future noise, but shuffled k-fold
can exploit overlap between adjacent future-horizon labels. Large positive
$R^2$ is therefore evidence of validation leakage, not model skill.

Table~\ref{tab:controlled} reports a Random Forest experiment with $n=1500$,
$H=20$, five outer folds, and seed 0. The overlap column is the mean fraction of
training rows whose label window overlaps any test label window, averaged across
folds.

\begin{table}[htbp]
\centering
\caption{Controlled leakage task. Positive $R^2$ is fabricated because the
target is unpredictable by construction.}
\label{tab:controlled}
\small
\begin{tabular}{p{0.19\textwidth}p{0.36\textwidth}rrr}
\toprule
Library & Splitter & Mean $R^2$ & Mean overlap & Folds \\
\midrule
\pkg{scikit-learn} & \code{KFold(shuffle=True)} & 0.918 & 1.000 & 5 \\
\pkg{scikit-learn} & \code{KFold(shuffle=False)} & -1.017 & 0.025 & 5 \\
\pkg{scikit-learn} & \code{TimeSeriesSplit} & -2.506 & 0.035 & 5 \\
\pkg{scikit-learn} & \code{TimeSeriesSplit(gap=20)} & -1.430 & 0.000 & 5 \\
\pcv{} & \code{PurgedKFold} & -0.870 & 0.000 & 5 \\
\pcv{} & \code{WalkForwardSplit} & -1.899 & 0.000 & 5 \\
\pcv{} & \code{CombinatorialPurgedCV} & -1.471 & 0.000 & 15 \\
\pkg{tscv} & \code{GapKFold(gap_before=20, gap_after=20)} & -1.217 & 0.000 & 5 \\
\pkg{timeseriescv} & \code{CombPurgedKFoldCV} & -0.894 & 0.004 & 15 \\
\pkg{timeseriescv} & \code{PurgedWalkForwardCV} & -1.543 & 0.000 & 4 \\
\bottomrule
\end{tabular}
\end{table}

The shuffled k-fold score of 0.918 is not a small optimism effect. It is a
complete failure of the validation design. The blocked and chronological
baselines remove most of the effect but still admit small amounts of overlap
unless a suitable gap is supplied. A fixed \code{TimeSeriesSplit} gap can solve
this particular constant-horizon toy problem, but it does not provide
label-aware intervals, variable horizons, group-purged folds, diagnostics, or
CPCV paths. The \pcv{} splitters remove the overlap by construction and return
negative $R^2$, which is the expected outcome for an unpredictable target
evaluated out of sample.

\subsection{A real smart-meter benchmark separates temporal leakage from entity generalization}

The second experiment uses the Low Carbon London smart-meter data set from UK
Power Networks and the London Datastore \citep{ukpn_lcl}. The prediction task is
half-hourly household electricity demand forecasting. Features include calendar
and lagged-load information, and the target is a forward-horizon mean. The
validation schemes compare pooled shuffled k-fold, blocked k-fold, walk-forward
validation, and held-out-household validation.

The full-population benchmark scans 167,932,474 raw rows, identifies 4,284
eligible Standard-tariff households with at least one year of data, draws 20
seeded subsamples of 60 households, and evaluates each validation scheme with
the same modeling harness. Table~\ref{tab:lcl} reports mean WAPE and 95\%
t-intervals. WAPE is $\sum |\hat{y}-y| / \sum |y|$, reported in percent.

\begin{table}[htbp]
\centering
\caption{Low Carbon London benchmark over 20 seeded subsamples of 60
households. Lower WAPE is better.}
\label{tab:lcl}
\begin{tabular}{lrrr}
\toprule
Metric & Mean & 95\% CI low & 95\% CI high \\
\midrule
Naive shuffled k-fold WAPE & 41.68 & 40.37 & 42.99 \\
Blocked k-fold WAPE & 42.43 & 41.07 & 43.80 \\
\code{WalkForwardSplit} WAPE & 42.36 & 41.01 & 43.71 \\
\code{GroupKFold} household WAPE & 44.92 & 43.38 & 46.45 \\
Temporal gap, WAPE points & 0.68 & 0.53 & 0.83 \\
Temporal gap, relative percent & 1.60 & 1.27 & 1.94 \\
Household gap, WAPE points & 2.56 & 2.08 & 3.03 \\
Household gap, relative percent & 6.03 & 4.93 & 7.12 \\
\bottomrule
\end{tabular}
\end{table}

By design, the result is less dramatic than the synthetic example. The temporal
leakage gap between shuffled k-fold and walk-forward validation is measurable
but small: 0.68 WAPE points, or 1.60\% relative to walk-forward WAPE. The
larger effect is the household gap. Scoring on unseen households is 2.56 WAPE
points worse than the pooled temporal estimate, or 6.03\% relative. This is the
more important conclusion for deployment: if the model will be used for
customers not seen during training, a purely temporal split answers a different
question.

\subsection{Cross-domain examples show that leakage-aware validation can also return negative results}

The repository also contains notebooks that exercise the same validation logic
on other public data sets. Table~\ref{tab:examples} summarizes the role of each
example. Some are designed to expose a large leakage effect; others show that a
leakage-aware split can correctly report a small or absent gap. The
``0.83--0.91'' range in the first row refers to the companion notebook's two
models, k-nearest neighbors and Random Forest, rather than to multiple random
seeds; the Random Forest-only benchmark in Table~\ref{tab:controlled} reports
0.918.

\begin{table}[htbp]
\centering
\caption{Reproducible examples included with the package.}
\label{tab:examples}
\small
\begin{tabular}{p{0.20\textwidth}p{0.24\textwidth}p{0.48\textwidth}}
\toprule
Example & Data source & Main validation lesson \\
\midrule
Synthetic leakage proof & Generated & k-nearest neighbors and Random Forest report $R^2$ of 0.83--0.91 on noise \\
Air quality & UCI air-quality data & A clock feature plus overlapping labels fabricates $R^2$ near 0.99 \\
Earthquakes & USGS catalogue & Magnitude history has no skill; purged splits reject the illusion \\
Smart meters & Low Carbon London & Household generalization dominates temporal leakage \\
Clinical mortality & PhysioNet Challenge 2012 & Whole-patient group holds are needed for patient-level inference \\
Predictive maintenance & NASA C-MAPSS & Walk-forward validation matches run-to-failure deployment \\
Rainfall & NOAA GHCN-Daily & One-day-ahead labels need purge and embargo buffers \\
Electricity load & PJM hourly load & CPCV paths expose score dispersion across backtest paths \\
Model comparison & Binance public bars & DSR prevents selecting an apparent edge after multiple trials \\
Sports prediction & Premier League matches & Honest validation shows calibration drift rather than a headline gap \\
\bottomrule
\end{tabular}
\end{table}

The examples deliberately include negative results. In the model-comparison
notebook, several models are tried on the same public price data. Once the
Deflated Sharpe Ratio corrects for the number of trials, no model clears a DSR
threshold of 0.95. In the PJM electricity-load notebook, CPCV produces five
paths whose DSR values range from 0.0011 to 0.7761 after correction for 20
trials. These are not failures of the software. They show that the validation
pipeline is capable of reporting that no reliable edge survived.

\section{Discussion}

The experiments show that leakage-aware validation is not a single recipe. In
the controlled task, randomization is catastrophic because every test label has
overlapping training labels. In the smart-meter benchmark, the temporal effect
is small but statistically visible, while the larger operational issue is
whether the model is expected to generalize to new households. In other domains,
the required split can be driven by patients, engines, seasons, stations, or
market regimes. The validation object should encode that deployment question
rather than being chosen only for convenience.

\pcv{} therefore treats diagnostics as first-class objects. A user can
construct a split with this package, with another package, or by hand, and then
check the interval and group invariants directly. This matters for
reproducibility. A reported model score is only as meaningful as the split that
created it, and the split should be auditable from code rather than described
informally in prose.

Several packages overlap with part of this problem. \pkg{scikit-learn} provides
\code{TimeSeriesSplit}; its \code{gap} argument is a fixed integer count rather
than a label-aware interval, and it does not provide group-purged folds, CPCV
paths, or split-level diagnostics. \pkg{tscv} provides fixed-gap splits, which
are useful when the required buffer is known and constant, but it does not
represent variable label horizons or grouped deployment targets.
\pkg{timeseriescv} implements purged and combinatorial time-series
cross-validation, but it does not unify variable-horizon label intervals,
group-purged folds, post-test embargoes, CPCV path reconstruction, and
independent diagnostic assertions in a typed \pkg{scikit-learn}-compatible
package \citep{timeseriescv2018}. \pkg{mlfinlab} is the best-known
implementation associated with the financial machine-learning literature, but
it is distributed as a commercial product and therefore cannot serve as a
permissive dependency for open scientific software \citep{mlfinlab2026}. The
companion benchmark also records two non-tabulated open alternatives:
\pkg{mlfinpy} did not run on the modern \pkg{pandas} stack used here, and
\pkg{RiskLabAI} failed because a plotting dependency was unavailable. Those
failures are recorded with exact exception messages rather than imputed scores.

\pcv{} is therefore not differentiated by claiming new purging mathematics. Its
contribution is integration and auditability. Unlike fixed-gap splitters or
single-purpose CPCV implementations, \pcv{} unifies variable-horizon label
intervals, group-purged folds, post-test embargoes, CPCV path reconstruction,
and split-level diagnostics as assertions that can be run on third-party or
hand-written splits. This combination is what lets the same validation contract
be used in ordinary \pkg{scikit-learn} model selection, in notebook examples,
and in automated tests.

There are limitations. Purging and embargoing remove a specific class of
validation leakage; they do not solve all forms of leakage. Feature engineering
can still use future data, target transformations can still be computed
globally, preprocessing can still be fit outside the training fold, and entity
leakage can still occur if the wrong group identifier is supplied. The package
does not claim that every chronological split is optimal. In highly
non-stationary settings, any historical validation estimate can be unstable. The
role of the package is narrower: when labels are interval-valued, it makes the
no-overlap condition explicit and executable.

Another limitation is maturity. The package is new, even though the underlying
methods are established. The open repository contains tests, type checks,
documentation, notebooks, and reproducible benchmarks, but wider external use
will be needed to discover edge cases in unfamiliar data layouts. For this
reason the software should be treated as validation infrastructure whose outputs
remain the analyst's responsibility, not as an automatic guarantee of scientific
validity.

Overlapping-label prediction problems require more than a chronological split.
The validation design must remove training labels that overlap test labels,
respect any post-test dependence buffer, and match the entity structure of the
deployment target. The practical contribution is not that the package makes
models better, but that it makes their validation harder to fool.

\section{Methods}

\subsection{Purging, embargoing, and group separation}

The implementation follows the half-open interval convention $I_i = [p_i,e_i)$.
A training row is purged from a fold if its label interval intersects any test
label interval. Boundary-touching intervals are not purged: if $e_i=p_j$ or
$e_j=p_i$, the intersection is empty under the half-open convention. This
matches common time-series indexing practice and prevents unnecessary deletion
when one forecast horizon ends exactly at the next prediction time.

Embargoing is applied after purging. It removes training rows whose prediction
times fall inside a post-test buffer. The buffer can be expressed as a duration
for timestamped data or as a fraction of the sample. In grouped validation,
group-purged splitters additionally require disjoint train and test groups, so
that deployment to unseen entities can be evaluated directly.

\subsection{Combinatorial purged cross-validation}

CPCV divides the ordered sample into $N$ contiguous blocks and holds out $k$
blocks at a time, producing $\binom{N}{k}$ test-block combinations. Each
combination defines a fold whose training side is purged and embargoed against
the union of its held-out test label intervals. Because the held-out blocks can
be non-adjacent, purging is performed against merged test intervals rather than
against the convex hull spanning all held-out blocks. The resulting fold
predictions can be recombined into several out-of-sample paths, giving a
distribution of validation trajectories rather than a single historical path.

\subsection{Implementation and diagnostics}

The package separates low-level interval operations from splitter orchestration.
The same interval functions are used by \code{PurgedKFold},
\code{PurgedGroupKFold}, \code{WalkForwardSplit}, and
\code{CombinatorialPurgedCV}. Diagnostic assertions accept train indices, test
indices, prediction times, evaluation times, and optional group labels. They can
therefore audit splits produced by \pcv{}, by another library, or by custom
code.

The package is maintained as a public open-source project with continuous
integration, strict static typing, and an extensive test suite covering split
invariants, numerical metrics, end-to-end reproducibility, notebook-derived
fixtures, and packaging quality gates. The repository accepts issues and pull
requests, and the core validation behavior is protected by tests rather than by
example output alone.

\subsection{Controlled leakage benchmark}

The controlled benchmark uses $n=1500$ observations and a horizon of $H=20$.
For each row $t$, the response is the mean of the next $H$ independent noise
values. The feature matrix contains only a monotone time index. Random Forest
regression is evaluated with five outer folds for ordinary k-fold,
time-series splits, fixed-gap baselines, and \pcv{} splitters. The reported
overlap metric is the mean fraction of training rows whose label window
overlaps any test label window. Competitor rows are generated by the repository
script \code{tools/competitor_benchmark.py}; unavailable competitors are
recorded as \code{NOT RUN} with exact exception messages.

\subsection{Low Carbon London benchmark}

The Low Carbon London benchmark uses the public smart-meter corpus released by
UK Power Networks and the London Datastore \citep{ukpn_lcl}. The benchmark
script scans the raw CSV files, retains Standard-tariff households with at
least one year of data, and samples 20 independent groups of 60 households. For
each subsample, the same feature builder, estimator, and scoring code are used
across pooled shuffled k-fold, blocked k-fold, walk-forward validation, and
held-out-household validation. The primary score is WAPE,
$\sum |\hat{y}-y| / \sum |y|$, reported in percent. Confidence intervals are
ordinary 95\% t-intervals across the 20 seeded subsamples.

\subsection{Cross-domain notebooks and statistical metrics}

The companion notebooks apply the same validation checks to public data from
earthquake catalogues, air-quality monitoring, smart meters, clinical mortality
prediction, engine run-to-failure simulation, rainfall records, electricity
load, public market bars, and sports prediction. The financial-statistical
metrics in \pcv{} implement the Probabilistic Sharpe Ratio, Deflated Sharpe
Ratio, and Minimum Track Record Length following Bailey, Lopez de Prado, and
coauthors \citep{bailey2012psr,bailey2014dsr}. These metrics are included
because model selection after many trials can itself create optimistic evidence,
even when a split is temporally valid.

\subsection{Computational environment and reproducibility}

The benchmark tables reported here were produced with \pcv{} 0.0.6, Python
3.12.7, \pkg{numpy} 2.4.5, \pkg{pandas} 3.0.3, \pkg{scikit-learn} 1.8.0, and
\pkg{scipy} 1.17.1 on macOS 26.3.1 (\code{arm64}). The current released
version of \pcv{} at the time of this manuscript is 0.0.10; the change set from
0.0.6 to 0.0.10 is limited to input-validation hardening, CI tooling, and
documentation, with no changes to the cross-validation algorithms or to the
numerical paths exercised by the reported computations. The package supports
Python 3.10 through 3.14; runtime dependency lower bounds are \pkg{numpy} 1.24,
\pkg{pandas} 2.0, \pkg{scikit-learn} 1.3, and \pkg{scipy} 1.10.

A split-generation microbenchmark is tracked as \code{tools/microbench.py}. It
uses 1,000,000 timestamped rows, five folds, one feature, a constant 20-second
label horizon, and no estimator fitting. In the recorded local run,
\code{PurgedKFold} generated the five folds in 1.898 seconds best-of-three
(mean 1.911 seconds), and wrote the environment details to
\path{paper/microbench_summary.md}. The full Low Carbon London benchmark
scanned 167,932,474 raw rows and ran 20 seeded 60-household subsamples in 53.8
minutes on the author's local machine.

The main local reproduction commands are:

\begin{Code}
pip install -e ".[dev,examples]"
pip install tscv timeseriescv
pytest -q
python tools/microbench.py
python tools/competitor_benchmark.py --out-dir examples/data
python tools/lcl_full_benchmark.py --k 20 --n 60 --seed 0
\end{Code}

The competitor command above reproduces the reported Table~\ref{tab:controlled}
rows when \pkg{tscv} and \pkg{timeseriescv} are installed. For a fast core-only
smoke check, use \code{--core-only}. The Low Carbon London command expects the
raw CSV files to be present locally. For faster checks, the repository includes
end-to-end tests with synthetic fixtures that exercise the same parser, feature
builder, and benchmark output format.

\subsection{Use of generative AI tools}

Large language model tools, including OpenAI Codex and ChatGPT in the GPT-5
family and Anthropic Claude in the Claude 4 family, were used as assistants for
code review, refactoring suggestions, test scaffolding, documentation drafting,
copy-editing, and pre-submission checks. These tools were not treated as
authors. All design decisions, AI-assisted changes, and outputs were reviewed
and validated by the author through unit, property, doctest, end-to-end,
type-checking, linting, and benchmark tests; no claim was accepted without
source or executable verification.

\section*{Data availability}

The synthetic benchmark data are generated by scripts included in the public
repository. The Low Carbon London raw smart-meter files are publicly available
from UK Power Networks and the London Datastore; they are not redistributed in
this repository because of their size. The repository contains the benchmark
scripts, derived summaries, and notebook fixtures needed to reproduce the
reported tables once the raw public files are available locally. The other
examples use public data sources identified in the notebooks and reference
list, including the U.S. Geological Survey catalogue \citep{usgs_catalog}, UCI
air-quality data \citep{devito2008airquality}, NOAA GHCN-Daily
\citep{noaa_ghcnd}, NASA C-MAPSS \citep{nasa_cmapss}, PhysioNet
\citep{physionet2012goldberger,physionet2012silva}, PJM hourly load data
\citep{pjm_load}, and Binance public market data via \pkg{pricehub}
\citep{binance_pricehub}.

\section*{Code availability}

The software is available from the project repository at
\url{https://github.com/eslazarev/purged-cross-validation} and distributed on
PyPI as \pkg{purgedcv} at \url{https://pypi.org/project/purgedcv/}. The source
distribution contains the examples and benchmark tools; the wheel contains the
importable \pcv{} package. The software is MIT licensed and archived on Zenodo
with DOI \doi{10.5281/zenodo.20312695}. The current released version at the
time of this manuscript is 0.0.10.

\bibliographystyle{unsrtnat}
\bibliography{paper}

\section*{Acknowledgements}

The author thanks the maintainers of the open data sets used in the
reproducible examples and the maintainers of the scientific Python ecosystem.
The purging, embargoing, CPCV, PSR, DSR, and MinTRL methods implemented in
\pcv{} are due to Lopez de Prado, Bailey, and colleagues; any implementation
errors are the author's.

\section*{Author contributions}

E.L. conceived the study, implemented the software, designed and ran the
experiments, analyzed the results, wrote the manuscript, prepared the
reproducibility materials, and approved the final version.

\section*{Additional information}

\subsection*{Competing interests}

The author declares no competing interests.

\subsection*{Funding}

No external funding was received for this work.

\subsection*{Ethics declarations}

This study did not involve new experiments with human participants, human
tissue, live vertebrates, or higher invertebrates. The empirical examples use
publicly released or synthetic data; the Low Carbon London data are used in the
public anonymized form released by the original data providers.

\end{document}
